The {\tt IMPEDANCE} element combines the physics of a co-located {\tt
ZLONGIT} (longitudinal monopole impedance) and one or more {\tt
ZTRANSVERSE} elements (transverse dipole or quadrupole
impedances) into a single element.  In a case where all these components
are modeled, using a single \verb|IMPEDANCE| element reduces computational
effort by about 40\% compared to three separate elements (single \verb|ZLONGIT|
plus two \verb|ZTRANSVERSE| elements for dipole and quadrupole).

\subsection*{Channels}

Five wake channels are available; each is controlled by a pair of
column-name parameters that select the (real, imaginary) columns of
the SDDS impedance file ({\tt INPUTFILE}, with the frequency column
named by {\tt FREQCOLUMN}):
\begin{center}
\begin{tabular}{lll}
\hline
parameter pair               & physical meaning & impedance units \\\hline
{\tt ZL\_REAL,  ZL\_IMAG}   & longitudinal monopole, $\Delta E \propto Z_L \cdot I$  & $\Omega$         \\
{\tt ZXD\_REAL, ZXD\_IMAG} & transverse $x$ dipole, $\Delta x' \propto Z_{xD}\cdot \langle x \rangle I$ & $\Omega/{\rm m}$ \\
{\tt ZYD\_REAL, ZYD\_IMAG} & transverse $y$ dipole, $\Delta y' \propto Z_{yD}\cdot \langle y \rangle I$ & $\Omega/{\rm m}$ \\
{\tt ZXQ\_REAL, ZXQ\_IMAG} & transverse $x$ quadrupole, $\Delta x' \propto x\cdot Z_{xQ}\cdot I$    & $\Omega/{\rm m}^2$ \\
{\tt ZYQ\_REAL, ZYQ\_IMAG} & transverse $y$ quadrupole, $\Delta y' \propto y\cdot Z_{yQ}\cdot I$    & $\Omega/{\rm m}^2$ \\
\hline
\end{tabular}
\end{center}
Any column-name parameter left at its default ({\tt NULL}) disables
that wake; at least one wake must be enabled.  A single SDDS file
typically holds all five channels.  The sign and Fourier conventions
match those of {\tt ZLONGIT} and {\tt ZTRANSVERSE} (see also
\verb|wake2impedance| and \verb|trwake2impedance| in the elegant
distribution).

\subsection*{Scale factors}

Following the {\tt LRWAKE} convention, every channel has its own
scale parameter:
\begin{itemize}
\item {\tt FACTOR} multiplies every channel.
\item {\tt ZFACTOR} multiplies the longitudinal channel only.
\item {\tt XFACTOR} ({\tt YFACTOR}) multiplies the $x$ ($y$) dipole
  channel only.
\item {\tt QXFACTOR} ({\tt QYFACTOR}) multiplies the $x$ ($y$)
  quadrupole channel only.
\end{itemize}
The effective scale of each channel is {\tt FACTOR} $\times$ the
channel factor $\times$ pass-ramp factor.

\subsection*{Binning}

The impedance file determines both the number of bins and the bin
width: with $N_f$ points spaced by $\Delta f$, the element uses $N_b =
2(N_f-1)$ bins of width $1/(N_b\,\Delta f)$.  The {\tt N\_BINS} and
{\tt BIN\_SIZE} parameters from {\tt ZLONGIT}/{\tt ZTRANSVERSE} are
present for compatibility but are overridden by the file.  Bins are
centered as in {\tt ZTRANSVERSE} (bin centers at
$t_{\rm min}+i\,dt$); the {\tt AREA\_WEIGHT} option is honored on the
longitudinal current histogram only.

\subsection*{Notes}

\begin{enumerate}
\item The beam charge must be specified via a {\tt CHARGE} element
  somewhere upstream in the beamline.  Unlike {\tt ZLONGIT} and {\tt
  ZTRANSVERSE}, {\tt IMPEDANCE} has no deprecated charge parameter;
  omitting the {\tt CHARGE} element is an error.
\item Only file-based impedances are accepted; the broad-band
  resonator inputs of {\tt ZLONGIT} and {\tt ZTRANSVERSE} are not
  implemented.  A resonator can be pre-tabulated onto a fine
  frequency grid (with one of {\tt elegant}'s helper scripts) and
  loaded through {\tt INPUTFILE}.
\item Output of the per-channel wake potentials to the file named by
  {\tt WAKES} is available only in the serial version, as for {\tt
  ZLONGIT} and {\tt ZTRANSVERSE}.  When written, one $V_t$ column is
  emitted per enabled channel ({\tt VtZL}, {\tt VtZxD}, {\tt VtZyD},
  {\tt VtZxQ}, {\tt VtZyQ}).
\item The original {\tt ZLONGIT} and {\tt ZTRANSVERSE} elements are
  unchanged and remain available.
\item Bunched-mode application is available via {\tt
  BUNCHED\_BEAM\_MODE=1}; see
  Section~\ref{sect:bunchedBeams}.
\item Numerical equivalence with stacked {\tt ZLONGIT} + {\tt
  ZTRANSVERSE} is exact to round-off for the transverse channels.
  For the longitudinal channel, {\tt IMPEDANCE} uses the bin-center
  convention of {\tt ZTRANSVERSE} rather than the left-edge convention
  of {\tt ZLONGIT}; this yields physically equivalent results that
  differ at the half-bin sampling level (typically $\sim 10^{-6}$
  relative).
\end{enumerate}
